\scanpage{147}
permet de \unsure{nous|les|se} ramener à un objet $P$ de la forme $\underline h(X)$, \deleted{\unsure{donc|dans} $u$ de la forme $\underline h(\alpha)$, où $\alpha\in\operatorname{Fl}(X,X)$. Or comme} \inserted{$\operatorname{Fl}(P,P)\simeq\operatorname{Fl}(X,X)$, donc} $\operatorname{Fl}(\deleted{\unsure{X}}P,\deleted{\unsure{X}}P)$ est un \uline{anneau semi-simple} $U$, \deleted{Or}

\deleted{Pour un tel anneau, \unsure{dont} il est \unsure{opérant} sur un objet d'une catégorie} \deleted{Alors $P$ et $u$ se décomposent suivant les composantes simples de $U=\operatorname{Fl}(P,P)$ (i.e. de son centre) et \unsure{on est ramené au cas où les}}

Or lorsque un anneau de ce type opère sur un objet d'une catégorie abélienne, alors pour tout $u\in U$, on peut trouver \deleted{deux couples d'\unsure{idempotents} orthogonaux complémentaires $\pi,\pi'$ de $U$, \unsure{tels que} $(\pi\pi'=\pi'\pi=0,\ \pi+\pi'=1,\ \pi^2=\pi,\ \pi'^2=\pi')$ tels que et $\varpi,\varpi'$ tels que $u\pi'=0,\varpi'u=0$,} $v\in U$ tel que
\[
\begin{aligned}
 vu&\ \text{soit un }\deleted{\text{projecteur}}\inserted{\text{\unsure{idem.pot|idempot|idempotent}}}\ \pi\\
 uv&\ \text{soit}\ \rule{2.5cm}{0.35pt}\ \varpi
\end{aligned}
\]
\begin{tikzpicture}[baseline=(phrase.base)]
\node[anchor=west,inner sep=0pt] (phrase) at (0,0) {et on ait les relations};
\node[draw,ellipse,inner sep=2pt] (relations) at (-0.3,-1.1) {\inserted{$\begin{aligned}\pi v&=v,\ u\pi=u\\\varpi u&=u,\ v\varpi=v\end{aligned}$}};
\draw[->] (relations.north east) .. controls +(0.5,0.5) and +(0.1,-0.7) .. (phrase.east);
\end{tikzpicture}

\diagram{\begin{tikzcd}[ampersand replacement=\&,row sep=small]
 P\simeq\pi P\unsure{+|\oplus}\pi'P \arrow[dr] \&\\
 \& \pi P\arrow[d,"\sim"]\\
 Q\simeq\varpi Q+\varpi'Q \& \varpi Q\arrow[l]
\end{tikzcd}}

\diagram{\begin{tikzcd}[ampersand replacement=\&,column sep=small,row sep=small]
P\arrow[r,"u",very thick]\arrow[d] \& \deleted{\unsure{P}}Q\arrow[l,bend right=38,"v"']\\
P'\arrow[r,"\sim"] \& Q'\arrow[u]
\end{tikzcd}}
\[
\left\{\begin{aligned}
 vu&=\text{projecteur }\pi,\\
 uv&=\text{projecteur }\varpi,\\
 &\deleted{\unsure{\varpi u}}
\end{aligned}\right.
\]
de cette façon, on gagne donc.

\diagram{\begin{tikzcd}[ampersand replacement=\&,column sep=small,row sep=small]
 P'+P''\arrow[dr] \& \& \& Q'+Q''\arrow[dl]\\
 \& P'\arrow[ul,shift right=1.6pt]\arrow[ul,shift left=4pt]\arrow[r,shift right=2pt,very thick] \& Q'\arrow[l,shift right=2pt]\arrow[ur,shift right=2pt] \&
\end{tikzcd}}
Pour les conditions de chaînes, on utilise le fait que les $\operatorname{Fl}(X,X)$ sont de \deleted{degré} dim finie sur $k$. ....
% archive sheet number: 146
